\scanpage{29}
\mspage{13}

correspondantes de relèvements \inserted{de $\{1,\dot\tau\}$} sont resp. les relèv\unsure{ts} en $\tau_0$, en $\tau$, et en $\tau_\infty$ -- qui ne donnent rien de nouveau. Reste à considérer $\tau\dot\sigma_1$ (et ses conjugués \inserted{par $\rho$} $\tau\dot\sigma_\infty,\tau\unsure{\dot\sigma_0}$), on a $\dot\sigma_1(2)=1/2$, donc

\[
(\tau\dot\sigma_1)(2)=\frac12
\]

\[
\text{et}\qquad U^{\tau\dot\sigma_1}=\{z\in U\mid z\bar z=1\}=\mathbb U-\{1\},
\]

où $\mathbb U=\{z\in\mathbb C\mid |z|=1\}$. Cet ensemble contient $P$ et $\bar P$, et ne donne donc pas de nouveau relèvement p.r. à celui déjà donné, avec $\widetilde\sigma_1$ relevant $\dot\sigma_1\tau$.

Il faudrait quand-$=$ donner aussi une description de $\unsure{\Sigma|E}$ tt entier par gén. et relations, en partant d'un syst. de gén. de $\Gamma=\unsure{\mathfrak S_3|\mathcal S_3}\times\mathbb Z/2$, qu'on va relever de façon judicieuse. On prendra les \inserted{quatre} générateurs

\deleted{$\rho'=\rho\tau,\unsure{\tau}\widetilde{\dot\sigma}_0,\widetilde{\dot\sigma}_1,\widetilde{\dot\sigma}_\infty$, avec les relations}

\diagram{\deleted{$\begin{array}{ccc}\Vert&\Vert&\Vert\\\dot\sigma_0\tau&\dot\sigma_1\tau&\dot\sigma_\infty\tau\end{array}$}}

\deleted{(donc $\rho'{}^6=1,\rho'{}^3=\tau,\rho'{}^2=\rho^2=\rho^{-1}$)}

\[
\left\{\begin{aligned}
&\deleted{\widetilde{\dot\sigma}_1^2=\widetilde{\dot\sigma}_{\unsure{2|1}}^2=\widetilde{\dot\sigma}_3^2=1}\\
&\deleted{\widetilde{\dot\sigma}_0\widetilde{\dot\sigma}_1=\widetilde{\dot\sigma}_1\widetilde{\dot\sigma}_\infty=\widetilde{\dot\sigma}_\infty\widetilde{\dot\sigma}_0=\dot\rho^{\unsure{-2|2}}}\\
&\deleted{\rho'{}^6=1}\\
&\deleted{\rho'{}^3\ \text{commute aux }\dot\sigma_i}
\end{aligned}\right.
\]
% archive sheet number: 28
